% TLDR: Define inputs, the human population, and ordinal forecasts without conflating proof validity.
\section{Target and Measurement Conditions}
\label{sec:target}
For task $t$ and solver configuration $a$, let
\begin{equation}
 z_{ta}=(x_t,\widehat y_{ta},y_t^*),
 \qquad R_{tah}\in\{0,1,2,3,4,5\}.
\end{equation}
Here $x_t$ is the original Python source and task contract, $\widehat y_{ta}$ is the candidate Lean artifact, $y_t^*$ is the reference, and $h$ identifies a human rater.
The normalized scalar rating is $R_{tah}/5$.
The main study displays the complete candidate and reference artifacts to raters and the judge, including any proof bodies or \texttt{sorry} placeholders.
Both receive identical input contents; the rendered view is recorded, and proof completion is a separate variable rather than the coverage target.

\subsection{What coverage means}
A formal specification is a theorem statement considered without its proof.
Coverage concerns the required statements and the domains and assumptions under which they hold.
A candidate may cover a reference statement with a stronger statement or with several statements together.
Counting theorem names or measuring textual overlap cannot establish this implication.
Conversely, narrowing a precondition can remove required cases even when a theorem's conclusion is unchanged.
Proof completion is recorded separately from statement coverage.

The source program constrains the intended behavior but does not automatically resolve a conflict between implementation behavior, task requirements, and the reference.
The protocol records such conflicts for independent inspection rather than silently choosing one input as infallible.
Conditioning on a reference makes the rating task more concrete; it can also make both humans and a judge prefer the reference's representation.
Reference sensitivity is therefore a measurement condition to test, not an assumed benefit of adding the reference.

\subsection{The human distribution}
Let $\mathcal H$ denote a specified population of qualified Lean and formal-methods experts using the same rubric.
The target is the probability mass function
\begin{equation}
 p^*(r\mid z)=\Pr_{h\sim\mathcal H}(R_h(z)=r).
\end{equation}
This definition requires a recruitment frame and a rating procedure, including allowed assistance, time guidance, and treatment of ambiguous inputs.
It does not make a convenience sample representative of all experts.
The empirical distribution from a fixed recruited panel is a different, narrower target; both must be labeled when reported.

The judge predicts $q_\theta(r\mid z)$, a nonnegative vector of six probabilities summing to one.
The initial interface elicits those six probabilities explicitly as structured output; it does not infer them from token likelihoods or the frequency of repeated scalar answers.
Any subsequently fitted probability mapping is trained on training ratings and selected on validation ratings, with its family and constraints recorded as a separate candidate condition.
Its expected normalized rating is $\mu_q(z)=\sum_{r=0}^5(r/5)q_\theta(r\mid z)$.
Each individual human rating is retained, including disagreement.
A histogram of stochastic scalar judge outputs does not by itself represent the human distribution: model-call instability and human disagreement have different sources.
The proposed probability-vector interface is consequently a new judge condition, not a reinterpretation of an old scalar score.
